\begin{statement}{G.23}{Proved}{Universal elliptic equations and exact finite CM transport leave the parameter free}
Let $k=\mathbb Q_5$, with the chosen embedding of $\mathbb Q(i)$,
and use $A_8,B_8$ from (G.24). For every $\lambda\in k$,
\[
             t\longmapsto\exp\bigl(t(A_8+\lambda B_8)\bigr)
\]
is a one-parameter subgroup of Enriquez's graded universal elliptic
Grothendieck--Teichm\"uller kernel $R_{\rm ell}^{\rm gr}$.
It commutes with $J(x)=ix$, $J(y)=-iy$; the interchange
$S(x)=y$, $S(y)=x$ acts on it by $t\mapsto-t$.
The full diagonal CM torus normalizes it through
$(a,b)\mapsto a^4b^4$.

Thus all defining universal elliptic group equations, together
with these character and finite-symmetry conditions, admit every
$\lambda$ to all orders. The same assertion about the leading
degree-eight vector holds for any prescribed lift of the finite
group $\langle J\rangle$, or $\langle J,S\rangle$, in the
corresponding $\mathrm{GL}_2$-extension of the positive
prounipotent kernel. The determinant-$-1$ element $S$ is taken
in this extension, not in the $\mathrm{SL}_2$ quotient of
$R_{\rm ell}^{\rm gr}$. These are universal group
elements, not asserted elements of the actual arithmetic subgroup.
\end{statement}
\begin{proof}
In \citet[Proposition 4.19, equation (48)]{Enriquez2014},
$\delta_{2n}$ is our $\epsilon_{2n+2}$: pairing the two halves
of his sum for its value on $y$ gives G.8. His bracket is
written in the opposite derivation algebra; changing its sign
does not change the generated Lie subalgebra or the membership
argument. The same section defines
$\mathfrak b_3=\langle\mathfrak{sl}_2,\delta_{2n}\rangle$ and
gives the subgroup inclusion
\[
 \exp(\widehat{\mathfrak b}_3^+)\rtimes\mathrm{SL}_2
                      \subset R_{\rm ell}^{\rm gr}.
\]
Both $A_8$, an iterated $\mathfrak{sl}_2$ bracket with
$\epsilon_8$, and $B_8$, a bracket of two members of the
$\epsilon_4$ module, belong to $\mathfrak b_3$.
Exponentiation converges in the positive-degree filtration.
This is actual all-order group membership, not only a tangent
vector solving finitely many equations. The elliptic-associator
isomorphisms of \citet[\S4.8 and Proposition 4.20]{Enriquez2014}
also transport these elements to the ungraded universal group,
with the canonical associated graded maps. No generation
conjecture for $\mathfrak b_3$ is used.

Both vectors have character $(4,4)$ and change sign under $S$.
Their exact associative expansions are checked in
\texttt{g\_degree8\_nonmet.py}. Hence $J$ fixes them and the
diagonal torus scales them by $a^4b^4$. The homogeneous
$\mathrm{SL}_2$-stable algebra $\widehat{\mathfrak b}_3^+$
is also $\mathrm{GL}_2$-stable, so these determinant-$-1$
and diagonal actions are defined in
$\exp(\widehat{\mathfrak b}_3^+)\rtimes\mathrm{GL}_2$.
Exponentiation proves the asserted identities.

Here is the finite-lift assertion with its convergence argument.
Let $H$ be either finite group just specified and $U$ the
positive-degree prounipotent kernel. Every section
$H\to U\rtimes H$ is conjugate by $U(k)$ to the linear section.
On a successive central vector layer, a cocycle satisfies
$a_{gh}=a_g+g a_h$. For $b=|H|^{-1}\sum_h a_h$ one has
$a_g=b-gb$, so a conjugation removes that layer.
Repeating in increasing degrees gives corrections tending to
the identity, whose product converges in $U(k)$ and conjugates
the entire section. Conjugating the one-parameter subgroups by
the same element supplies the prescribed finite action.
Their leading degree-eight terms remain unchanged, because
the conjugator has positive degree.

The analytic elliptic associator at the CM point, its comparison
path with the Tate tangent, and its arithmetic realization carry
additional information. This proof does not specify those data.
It rules out selecting $\lambda$ from the stated universal
algebraic conditions, even with exact finite CM lifts, and
does not assert arithmetic realization of the candidate family.
\end{proof}

\begin{statement}{G.24}{Proved}{A saturated ambient plane and two finite nonmetabelian coordinates}
Let $[w]P$ denote an associative-word coefficient, and put
\[
 \begin{aligned}
 w_0&=x^5y^4,& w_1&=x^4yxy^3,\\
 \alpha(D)&=-[w_0]D(x),&
 \nu(D)&=[w_1]D(x)+4[w_0]D(x),\\
 C_8&=A_8+\frac{31}{20}B_8.
 \end{aligned}
\]
For $H=\mathbb Z_5x\oplus\mathbb Z_5y$, the plane
$W_8=\mathbb Q_5A_8\oplus\mathbb Q_5B_8$ has saturated ambient lattice
\[
 W_8\cap\operatorname{Hom}_{\mathbb Z_5}(H,L_9H)
               =\mathbb Z_5C_8\oplus\mathbb Z_5B_8.
\]
The coordinate matrix is
\[
 \begin{array}{c|rrr}
       &A_8&B_8&C_8\\\hline
 \alpha&1&0&1\\
 \nu&-31/20&1&0
 \end{array}.
\]
For an actual Galois element $\sigma$ in the eighth Johnson
subgroup these functionals isolate its character-$(4,4)$ part,
and
\[
 \begin{aligned}
 \alpha(D_8(\sigma))&=\frac{d^3\chi^{(4)}_4(\sigma)}{267840},\\
 \nu(D_8(\sigma))&=\alpha(D_8(\sigma))
                         \left(\lambda-\frac{31}{20}\right).
 \end{aligned}
\]
Thus any such element with nonzero $\alpha$ would determine
$\lambda=31/20+\nu/\alpha$.

For each $r\ge1$, both coordinates modulo $5^r$ are read from
the action on a characteristic finite quotient of $P_5$, namely
its image in the units of
\[
 (\mathbb Z/5^r)\langle X,Y\rangle/(X,Y)^{10},
 \qquad a\longmapsto1+X,\quad b\longmapsto1+Y,
\]
where $a,b$ lift the chosen integral homology basis.
This measurement uses no logarithm or division by $5$.
The lattice above is ambient; neither the arithmetic image
lattice nor the arithmetic value of $\nu$ is determined here.
\end{statement}
\begin{proof}
The exact tensor calculation gives
\[
 \begin{array}{c|rrr}
       &A_8(x)&B_8(x)&C_8(x)\\\hline
 [w_0]&-1&0&-1\\
 [w_1]&49/20&1&4
 \end{array}.
\]
Both $C_8$ and $B_8$ annihilate $[x,y]$. All their standard
Lyndon coefficients have denominators in $\{1,2,4\}$; the full
Lyndon vectors are in \texttt{code/g\_degree8\_nonmet.json}
and \texttt{code/g\_degree8\_results.json}. Hence both are integral
Lie derivations at $5$. The integral functionals $\alpha,\nu$
have the identity matrix on the pair $(C_8,B_8)$. Every integral
$D\in W_8$ therefore has
\[
                  D=\alpha(D)C_8+\nu(D)B_8
\]
with coefficients in $\mathbb Z_5$. This proves saturation;
it makes no assertion of arithmetic surjectivity.

Each selected word in $D(x)$ has five $x$'s and four $y$'s,
so both functionals kill all other homogeneous CM characters.
Formula (G.25) now gives the two arithmetic identities, with
its corrected theta-character normalization unchanged.
The second coordinate is nonmetabelian since $\nu(B_8)=1$
and the metabelian character kills $B_8$. It does not vanish
on our chosen lift $A_8$; the explicit $31/20$ corrects that
fact.

Choose lifts $a,b$ of $x,y$ in $P_5$. They form a basis modulo
the Frattini subgroup and are free pro-$5$ generators.
The completed group algebra identifies with the completed
noncommutative power-series algebra by $a\mapsto1+X$,
$b\mapsto1+Y$. The inverse maps follow directly from the free
pro-$5$ universal property and the substitutions
$X\mapsto a-1$, $Y\mapsto b-1$.
The ideal generated by $5^r$ and the tenth power of the
augmentation ideal is characteristic. Its quotient algebra
is finite, and the group $1+(X,Y)$ in it is a finite $5$-group:
its successive augmentation quotients are additive finite
$5$-groups. This proves the claimed characteristic finite
group quotient and the descent of the Galois action.

Since $\sigma$ is in the eighth Johnson subgroup,
$\sigma(a)a^{-1}\in\gamma_9P_5$. The length-nine part of its
Magnus series is $D_8(\sigma)(X)$; all positive lower parts
vanish. Indeed the integral associated graded of a free
pro-$5$ group is its free Lie algebra, and leading group
commutators expand as the corresponding tensor Lie brackets.
Applying $-[X^5Y^4]$ and $[X^4YXY^3]+4[X^5Y^4]$ to that
series gives the two coordinates modulo $5^r$.
Changing a logarithmic expansion cannot change this leading
Johnson term.

The code also verifies detection using actual group words.
At modulus $125$ and length nine, products of the standard
Lyndon group commutators with the coefficients of $C_8$ and
$C_8+B_8$ give central corrections to the two generators.
Their denominators are powers of $2$, so these powers of
commutators are well-defined. The program recursively expands
each group commutator, checks its leading Lie polynomial, and
checks that the resulting substitutions have inverses and
preserve the peripheral commutator. Their measured pairs
$(\alpha,\nu)$ are $(1,0)$ and $(1,1)$ modulo $125$.
These are ambient finite-quotient automorphisms, not claimed
Galois elements. G.15 ensures a nonzero arithmetic first
coordinate in the rational graded image, but its integral
image ideal and the corresponding second coordinate remain
unevaluated. The value of $\lambda$ therefore remains Open.
\end{proof}

